\chapter{Переход как первичный объект и научный язык композиции}

Предыдущие главы постепенно заменили образ частицы как хранимого вещества
образом устойчивого процесса. Настоящая глава формулирует эту замену
точнее: первичным кандидатом объявляется не состояние, не свет и не
планковская частица, а локальный допустимый переход. Наблюдаемое состояние
может возникать как угол или составной результат сети переходов, а частица
--- как устойчивый дефект закона их композиции.

Это пока не новая физическая теория. Задача главы --- определить, какими
уже существующими научными языками можно выразить такую онтологию и какие
из них действительно совпадают с построенными объектами проекта.

\section{Антикруговая граница}

В начале Тома V уже проверялась категория редукций между геометрическим,
спектральным и корреляционным чтениями. Минимальный треугольник не замкнулся:
спектральные данные не восстанавливают геометрию однозначно. Возвращение к
тому же треугольнику под новым названием было бы круговым движением.

Нынешняя постановка иная. Объектами не являются целые теории, а локальные
алгебры наблюдения $A_x$. Стрелкой является физически типизированное
соответствие
\begin{equation}
 {}_{A_y}E_{A_x},
 \label{eq:v5-transition-correspondence}
\end{equation}
а последовательные переходы составляются внутренним тензорным
произведением
\begin{equation}
 {}_{A_z}F_{A_y}\widehat\otimes_{A_y}{}_{A_y}E_{A_x}.
 \label{eq:v5-transition-composition}
\end{equation}
Такой язык стандартен для гильбертовых $C^*$-бимодулей и соответствий;
см. \path{arXiv:math-ph/0506024}. Категориальный язык квантовых процессов
также принимает системы за объекты, процессы за морфизмы и композицию за
основную операцию; см. \path{arXiv:quant-ph/0402130}. Но сама категория
задаёт грамматику, а не выбирает физическую динамику.

\section{Первый инструмент: соответствия Мориты}

Проект уже содержит конкретный переходный носитель
\begin{equation}
 E=M_{20\times15}(\mathbb C),
\end{equation}
соединяющий углы $M_{20}(\mathbb C)$ и $M_{15}(\mathbb C)$. Связывающая
алгебра имеет вид
\begin{equation}
 \mathcal L(E)=
 \begin{pmatrix}
 M_{20}(\mathbb C)&E\\
 E^*&M_{15}(\mathbb C)
 \end{pmatrix}.
 \label{eq:v5-transition-linking-algebra}
\end{equation}

Здесь внедиагональные элементы являются переходами. Их композиции
возвращаются в диагональные углы:
\begin{equation}
 EE^*=M_{20}(\mathbb C),
 \qquad
 E^*E=M_{15}(\mathbb C).
 \label{eq:v5-transitions-generate-corners}
\end{equation}
Следовательно, в строгом алгебраическом смысле наблюдаемые на углах могут
порождаться составлением переходов. Машинный аудит проверяет этот факт на
матричных единицах представителя $M_{3\times2}$: произведения двух
противоположно ориентированных стрелок линейно порождают все $9$ и все $4$
направления соответствующих углов.

Это наиболее прямой научный опорный пункт новой интуиции. Он не доказывает,
что мир ``состоит из переходов'', но показывает, что текущая архитектура
уже математически допускает чтение диагональных объектов как композитов
внедиагональных процессов.

\section{Второй инструмент: группоиды, колчаны и пути}

Для хранения информации о начале, конце и композиции маршрутов пригодны
группоиды и колчаны. Стрелка $e:x\to y$ имеет обратную стрелку, если
переход обратим, а замкнутый путь
\begin{equation}
 \gamma=e_n\cdots e_2e_1
\end{equation}
несёт голономию
\begin{equation}
 \Omega_\gamma=U_{e_n}\cdots U_{e_2}U_{e_1}.
 \label{eq:v5-transition-holonomy}
\end{equation}
Группоидная свёрточная алгебра исторически является естественной алгеброй
переходов; связь с матричной механикой и некоммутативной геометрией
обсуждается в \path{arXiv:math/0601054}.

При локальной смене базиса $g_x$ стрелка преобразуется как
\begin{equation}
 U_e\longmapsto g_yU_eg_x^{-1},
\end{equation}
а голономия замкнутого пути --- сопряжением. Поэтому её спектр, след и
класс сопряжённости не зависят от локального описания. Именно такие данные
могут наследоваться от фундаментальной сети наблюдаемым частицам.

\section{Третий инструмент: квантовые блуждания и клеточные автоматы}

Соответствие говорит, \emph{какие} переходы типизированы правильно, но не
говорит, \emph{когда} и \emph{как} они совершаются. Для этого нужен
локальный закон обновления. Обратимые квантовые клеточные автоматы задают
дискретный глобальный шаг из строго локальных унитарных правил с конечной
скоростью распространения; см. \path{arXiv:quant-ph/0405174}. Квантовое
блуждание является их одночастичным сектором.

В проекте этот инструмент уже прошёл первый тест. Оператор
\begin{equation}
 W(k)=
 \begin{pmatrix}
 ne^{ik}&-im\\
 -im&ne^{-ik}
 \end{pmatrix},
 \qquad m^2+n^2=1,
 \label{eq:v5-transition-primitive-walk}
\end{equation}
локален, унитарен и имеет дираковский предел
\begin{equation}
 E^2=p^2+M^2.
\end{equation}
Это означает, что язык локальных переходов способен породить привычную
релятивистскую кинематику. Но значение $m$ остаётся свободным. Поэтому
квантовый автомат является кандидатом на динамическую грамматику, а не
готовым объяснением масс.

\section{Четвёртый инструмент: индекс и устойчивый дефект}

Если частица является не отдельным переносимым веществом, а нарушением
однородного правила обновления, нужно различать случайную неоднородность и
устойчивый класс. Теория индекса квантовых блужданий и клеточных автоматов
измеряет чистый поток квантовой информации и классифицирует компоненты
локально обратимой динамики; см. \path{arXiv:0910.3675}.

Проектный кандидат на частицу должен поэтому задаваться не произвольным
профилем $m(x)$, а инвариантом, который нельзя удалить локальной сменой
базиса или непрерывной деформацией без закрытия щели. Голономия, индекс,
класс когомологии и знак Пфаффиана являются допустимыми типами такого
свидетельства. Локализованная волновая функция без защищающего инварианта
недостаточна.

\section{Четырёхтактное наследование}

Пусть фундаментальный шаг $U$ на внутреннем четырёхсостояний носителе
удовлетворяет $U^4=I$. Тогда характеры циклической группы дают четыре
канонических проектора
\begin{equation}
 P_\lambda=\frac14\sum_{n=0}^3\lambda^{-n}U^n,
 \qquad
 \lambda\in\{1,i,-1,-i\}.
 \label{eq:v5-transition-character-projectors}
\end{equation}
Они попарно ортогональны и суммируются в единицу. Поэтому ``наследование
подобия'' можно заменить точным выражением \emph{наследование характера
представления}: наблюдаемый сектор получает не геометрическую форму
первичного цикла, а его фазу, знак, ориентацию или индекс.

В частности, спинорный знак после двух шагов и полный возврат после
четырёх могут принадлежать разным проекторным чтениям одного $U$. Но
текущий проект ещё не вывел обязательное назначение этих четырёх
характеров фотону, фермионам или поколениям.

\section{Минимальный стек и отложенные инструменты}

Причинные множества дают язык, в котором порядок событий первичен, а
гладкое пространство должно возникнуть позднее; последовательная динамика
роста построена в \path{arXiv:gr-qc/9904062}. Спиновые пены, тензорные сети
и групповая теория поля также могут описывать сборку пространства из
локальных данных. Однако сейчас они добавили бы новые степени свободы до
того, как проверен уже имеющийся переходный носитель.

Поэтому минимальный стек Тома V замораживается так:
\begin{equation}
 \boxed{
 \begin{gathered}
 \text{соответствие Мориты --- тип стрелки},\\
 \text{колчан/группоид --- композиция маршрутов},\\
 \text{квантовый автомат --- локальный закон шага},\\
 \text{голономия и индекс --- идентичность дефекта}.
 \end{gathered}}
 \label{eq:v5-transition-language-stack}
\end{equation}
Категориальная теория процессов служит общей грамматикой этой композиции,
но не заменяет четыре конкретных слоя.

\section{Рабочая онтологическая формула}

Накопленные результаты допускают следующую, всё ещё условную формулу:
\begin{quote}
первичен локальный обратимый переход; свободное поле есть однородный режим
его повторения; частица есть устойчивый дефект композиции; наблюдаемые
квантовые числа являются характерами и индексами этого дефекта.
\end{quote}

Эта формула сильнее метафоры и слабее физической теории. Она уже указывает,
какие объекты должны быть построены, но пока не выводит локальную сеть,
масштаб шага, набор частиц или взаимодействия Стандартной модели.

\begin{theorem}[Научная достаточность языка переходов]
Существующие математические инструменты достаточны, чтобы строго задать
типизированные переходы, их композицию, локальную обратимую динамику и
топологически устойчивый дефект. Более того, все четыре слоя имеют
конкретные представители в текущем проекте. Однако они не фиксируют
единственный локальный закон обновления и не доказывают, что наблюдаемые
частицы исчерпываются дефектами такого закона.
\end{theorem}

Следующий конструктивный вопрос теперь однозначен: существует ли на
$E=M_{20\times15}(\mathbb C)$ минимальное нелинейное локальное правило,
которое само создаёт, переносит и сохраняет ненулевой индексный дефект,
не вводя профиль массы вручную? Если нет, онтология переходов останется
полезным языком организации, но не станет физическим источником материи.

\begin{protocol}
\begin{enumerate}
 \item отличие от закрытой категории редукций: \textbf{зафиксировано};
 \item тип перехода: \textbf{соответствие Мориты};
 \item композиция переходов: \textbf{внутренний тензорный продукт и
 связывающая алгебра};
 \item локальный закон времени: \textbf{унитарный квантовый автомат};
 \item устойчивое имя дефекта: \textbf{голономия или индекс};
 \item состояния как композиции переходов: \textbf{алгебраически
 реализовано на углах};
 \item четырёхтактное наследование: \textbf{реализовано проекторами
 характеров};
 \item единственный физический автомат: \textbf{не выведен};
 \item спектр частиц: \textbf{не выведен};
 \item физическое замыкание: \textbf{не достигнуто}.
\end{enumerate}
\end{protocol}

Машинный протокол сохранён в
\path{s2t_v5_transition_primitive_scientific_language_gate_results.json}.